\subsection{THETA3-C12: four non-CM factors and a period of 500}\label{C12}
Use Schoen's surface $Y$ and curves from pp.~133--134 of the cited paper. The two projections to the Fermat quintic $C_5$ are
$(y_{01}:y_{03}:y_{13})$ and $(y_{02}:y_{03}:y_{23})$.
Let $C_0$ be its degree-five quotient, with equation $y^2=1-x^5$, and let $C\to C_0$ be Schoen's etale double cover, obtained from the elliptic curve
\[
 F:\quad s^2=1+x+x^2+x^3+x^4.
\]
The source proves $j(F)=2048$ and that $F$ is non-CM. There is a degree-two map $C\to F$.

There is a connected etale degree-four cover $X\to Y$ and a generically finite map $(a_1,a_2):X\to F^2$ of degree~500. If
$\int_F i\alpha\wedge\bar\alpha=1$, then
\begin{equation}\label{eq:500}
 \int_X(a_1^*\alpha\wedge a_2^*\alpha)
 \wedge\overline{(a_1^*\alpha\wedge a_2^*\alpha)}=500.
\end{equation}
All four degree-one sources are non-CM. The map
$(a_1,a_2,a_1,a_2):X\to F^4$ pushes its fundamental cycle to
$500[\Delta_{13}\cap\Delta_{24}]$, a product of diagonal divisor classes.
\begin{proof}
On $y_{03}=1$ write $u=y_{13}$ and $v=y_{23}$. The function field of $C_5^2$ has equations
$y_{01}^5=1-u^5$, $y_{02}^5=1-v^5$. Adjoining $y_{12}$ with
$y_{12}^5=u^5-v^5$ gives the function field of $Y$. The radicand has valuation~1 at the generic point of the diagonal of $C_5^2$, since $u^5-v^5=(u-v)$ times a unit there. It is not a fifth power, so the Kummer extension has degree~5. Resolution leaves function fields unchanged. Hence the product of the two curve maps gives a morphism $Y\to C_0^2$ of generic degree $5\cdot5\cdot5=125$.

The connected Galois etale cover $C^2\to C_0^2$ has degree~4 and group $(\Z/2)^2$. Its function field is linearly disjoint from the degree-125 extension, since the intersection degree divides both~4 and~125. Therefore
$X=Y\times_{C_0^2}C^2$ is connected and etale of degree~4 over $Y$. It is smooth projective and remains an arithmetic ball quotient. Its generic degree over $C^2$ is~125, and the two degree-two maps to $F$ give generic degree $125\cdot2\cdot2=500$ over $F^2$.

For $\eta=\mathrm{pr}_1^*\alpha\wedge\mathrm{pr}_2^*\alpha$ the Koszul sign gives
$\int_{F^2}\eta\wedge\bar\eta=-(-i)^2=1$.
Pushforward of the fundamental class under a proper generically finite morphism multiplies it by its degree, proving \eqref{eq:500}. Postcomposing with $(p,q)\mapsto(p,q,p,q)$ proves the cycle formula. The image is the transverse intersection of the two diagonal divisors.
\end{proof}
\framework{This realizes a four-factor non-CM period by an algebraic surface correspondence, in the sense of cote~44, scans~7--8. Its class is already generated by divisors. The new forms arise on a noncongruence cover, rather than from the original adelic character lifts. Thus neither a new exceptional Hodge-class case nor an extension of the original congruence theta packet is asserted.}
